% answers_short.tex —— 2006 年·普通高等学校招生考试·辽宁卷理科 · 简短答案速查
% 由 tools/make_short_answers.py 自 answers.tex 的 \gksolhead/\gkanswer 提取，
% 并与 2006/2006_数学真题_辽宁理_答案要点.md 逐题比对后生成，勿手改（改源 MD 后重跑）。

\begin{gknotice}
\textbf{使用说明}：本节只给最终答案便于快速核对，完整推导见解析版。答案为编写者独立推导，\textbf{非官方标准答案}。
\end{gknotice}

\gkpart{选择题}
\begin{center}\begin{tabular}{ccccccccccccc}
\toprule
题号 & 1 & 2 & 3 & 4 & 5 & 6 & 7 & 8 & 9 & 10 & 11 & 12 \\
\midrule
答案 & C & D & D & A & C & B & A & A & C & D & C & B \\
\bottomrule
\end{tabular}\end{center}

\gkpart{填空题}
\begin{itemize}[leftmargin=2.6em]
  \item[\textbf{13.}] $1/2$。
  \item[\textbf{14.}] $-1$。
  \item[\textbf{15.}] 48。
  \item[\textbf{16.}] $\sqrt{2/3}$。
\end{itemize}

\gkpart{解答题}
\begin{itemize}[leftmargin=2.6em]
  \item[\textbf{17.}] (1) 最大值 $2+\sqrt2$，取得最大值时 $x=\pi/8+k\pi$，$k\in\mathbb Z$；(2) 单调增区间为 $[-3\pi/8+k\pi,\ \pi/8+k\pi]$，$k\in\mathbb Z$
  \item[\textbf{18.}] (1) $BF\parallel$ 平面 $ADE$；(2) $G\in EF$，$\cos\theta=1/4$
  \item[\textbf{19.}] (1) $\xi_1$ 取 1.2、1.18、1.17 的概率依次为 $1/6,1/2,1/3$；$E\xi_1=1.18$。$\xi_2$ 取 1.3、1.25、0.2 的概率依次为 $(1-p)^2,2p(1-p),p^2$；$E\xi_2=1.3-0.1p-p^2$ 万元；(2) $0<p<0.3$
  \item[\textbf{20.}] (1) 线段 $AB$ 是圆 C 的直径；(2) $p=2$
  \item[\textbf{21.}] (1) $x_0=-1$；(2) $a=3\sqrt3,\ d=3$
  \item[\textbf{22.}] (1) $f_k(1)=1$；(2) 对题定范围内任意 $x_1,x_2\in[-1,1]$，$|F(x_1)-F(x_2)|\le 2^{n-1}(n+2)-n-1$
\end{itemize}
